\documentclass[a4paper,12pt]{article}
\usepackage[utf8]{inputenc}
\usepackage[ngerman]{babel}
\usepackage{amsmath}
\usepackage{parskip}

\setlength{\parindent}{0pt}

\title{Modul 294 - Aufgabenblatt 14}
\author{von Lukas Meier}
\date{Unterricht vom 31.03.2026}

\begin{document}

\maketitle

\section{Vorbereitung}

Erstellen Sie einen neuen Projektordner fuer diese Lektion.

Legen Sie mindestens folgende Struktur an:
\begin{verbatim}
index.html
js/
  main.js
\end{verbatim}

Verknuepfen Sie \texttt{main.js} im HTML mit:

\begin{verbatim}
<script type="module" src="./js/main.js"></script>
\end{verbatim}

\section{Warum modularisieren?}

Beantworten Sie kurz in eigenen Worten:
\begin{itemize}
    \item Warum ist eine einzelne sehr grosse \texttt{main.js} problematisch?
    \item Welche Vorteile bringen mehrere kleine Module?
\end{itemize}

\section{Erstes Modul erstellen}

Erstellen Sie die Datei \texttt{js/math.js}.

Exportieren Sie darin die Funktionen:
\begin{itemize}
    \item \texttt{add(a, b)}
    \item \texttt{subtract(a, b)}
    \item \texttt{multiply(a, b)}
\end{itemize}

Importieren Sie diese Funktionen in \texttt{main.js} und geben Sie mehrere Testergebnisse in der Konsole aus.

\section{Named Exports verwenden}

Stellen Sie sicher, dass Sie fuer die Funktionen aus \texttt{math.js} Named Exports verwenden.

Testen Sie zusaetzlich einen Import mit Alias, zum Beispiel:
\begin{itemize}
    \item \texttt{add as plus}
\end{itemize}

\section{Default Export erstellen}

Erstellen Sie die Datei \texttt{js/formatPrice.js}.

Exportieren Sie eine Funktion \texttt{formatPrice(value)} als Default Export.

Die Funktion soll Zahlen in folgendes Format umwandeln:
\begin{itemize}
    \item \texttt{12 -> CHF 12.00}
    \item \texttt{3.5 -> CHF 3.50}
\end{itemize}

Importieren Sie die Funktion in \texttt{main.js} und testen Sie sie.

\section{Code in Partials aufteilen}

Erstellen Sie fuer eine kleine Produktliste zusaetzliche Dateien:
\begin{itemize}
    \item \texttt{js/data/products.js}
    \item \texttt{js/ui/renderProducts.js}
    \item \texttt{js/utils/calculateTotal.js}
\end{itemize}

Inhalt der Module:
\begin{itemize}
    \item \texttt{products.js}: exportiert ein Array mit mindestens drei Produkten
    \item \texttt{renderProducts.js}: enthaelt eine Funktion zum Rendern der Produkte ins HTML
    \item \texttt{calculateTotal.js}: berechnet die Summe aller Produktpreise
\end{itemize}

\section{HTML vorbereiten}

Erweitern Sie \texttt{index.html} um folgende Elemente:
\begin{itemize}
    \item eine Ueberschrift
    \item ein leeres \texttt{<ul>} oder \texttt{<div>} fuer die Produktliste
    \item ein Element fuer die Gesamtsumme
\end{itemize}

Vergeben Sie sinnvolle IDs, damit die Elemente per JavaScript angesprochen werden koennen.

\section{Module zusammenfuehren}

Nutzen Sie \texttt{main.js} als Einstiegspunkt.

Importieren Sie dort:
\begin{itemize}
    \item die Produktdaten
    \item die Render-Funktion
    \item die Funktion fuer die Gesamtsumme
    \item den Default Export aus \texttt{formatPrice.js}
\end{itemize}

Beim Laden der Seite soll:
\begin{itemize}
    \item die Produktliste angezeigt werden
    \item die Gesamtsumme berechnet werden
    \item die Gesamtsumme formatiert ins HTML geschrieben werden
\end{itemize}

\section{Reflexion zur Struktur}

Begruenden Sie kurz Ihre Dateiaufteilung:
\begin{itemize}
    \item Welche Datei enthaelt Daten?
    \item Welche Datei enthaelt Logik?
    \item Welche Datei ist fuer die Darstellung zustaendig?
\end{itemize}

\section{Fehlersuche}

Pruefen Sie Ihr Projekt gezielt auf folgende Punkte:
\begin{itemize}
    \item Ist im HTML wirklich \texttt{type="module"} gesetzt?
    \item Stimmen alle relativen Pfade?
    \item Haben alle Imports die Dateiendung \texttt{.js}?
    \item Stimmen Named Export und Named Import exakt ueberein?
\end{itemize}

Notieren Sie mindestens einen Fehler, den Sie waehrend der Arbeit gemacht und behoben haben.

\section{Abschluss}

Die Lektion ist abgeschlossen, wenn:
\begin{itemize}
    \item die App aus mehreren JS-Dateien besteht
    \item Imports und Exports korrekt funktionieren
    \item die Produktliste im Browser erscheint
    \item die Gesamtsumme korrekt berechnet und formatiert wird
\end{itemize}

\section{Optional}

Falls Zeit bleibt:
\begin{itemize}
    \item Erstellen Sie ein weiteres Modul \texttt{filterProducts.js}
    \item Filtern Sie alle Produkte ueber einem bestimmten Preis
    \item Teilen Sie DOM-Selektoren ebenfalls in ein eigenes Modul aus
\end{itemize}

\end{document}
